\chapter*{Summary}

In the previous phase, the initial Adaptive AUTOSAR project was developed and evaluated on a standard host machine via Docker, utilizing SOME/IP for communication. Building upon that foundation, this report presents the outcomes of the final phase of the Computer Engineering Capstone Project. This new phase significantly upgrades the automated driving software system by successfully deploying the entire architecture directly onto a resource-constrained edge device. Furthermore, it transitions to Data Distribution Service (DDS) middleware, deploys highly optimized, lightweight AI models that perform substantially faster than those in Phase 1, and integrates broader Intelligent Transportation Systems (ITS) information. The primary goal remains providing a clear, step-by-step guide on structuring advanced Adaptive AUTOSAR systems using an industrial toolchain.

To achieve this, the report details two major technical contributions. First, it provides a comprehensive breakdown of the configuration workflow utilizing the PopcornSAR toolchain, demonstrating how to establish an Adaptive AUTOSAR system from defining the high-level architecture to generating deployment artifacts. Second, it focuses on customizing lightweight AI models tailored specifically for embedded edge deployment. This ensures the perception system operates efficiently on physical hardware without compromising its seamless integration into the automotive software framework.

Finally, a functional prototype is introduced to demonstrate the practical viability of these methods on the edge device. By establishing a robust service-oriented communication link between a provider application and a client node within the Adaptive AUTOSAR ecosystem, the prototype serves as a critical proof of concept. Ultimately, this report establishes a rigorous methodological baseline for future intelligent vehicle developments, successfully merging configuration-driven AUTOSAR practices with optimized edge AI deployment.

\clearpage
